% SPDX-License-Identifier: MIT-0
\documentclass[11pt]{article}
\usepackage[T1]{fontenc}
\usepackage[utf8]{inputenc}
\usepackage{lmodern}
\usepackage[a4paper,margin=2.5cm,headheight=14pt]{geometry}
\usepackage{amsmath,amssymb,amsthm,mathtools}
\usepackage{microtype,booktabs,array,longtable}
\usepackage[dvipsnames]{xcolor}
\usepackage{titlesec,fancyhdr,enumitem}
\usepackage[colorlinks=true,linkcolor=MidnightBlue,citecolor=MidnightBlue,
 urlcolor=MidnightBlue,bookmarksnumbered=true]{hyperref}
\definecolor{Ink}{HTML}{16324F}
\definecolor{Accent}{HTML}{4E5C3D}
\titleformat{\section}{\Large\bfseries\color{Ink}}{\thesection}{0.7em}{}
\titleformat{\subsection}{\large\bfseries\color{Accent}}{\thesubsection}{0.7em}{}
\pagestyle{fancy}
\fancyhf{}
\fancyhead[L]{\small\color{Ink}Shellings of preorder lattice-point posets}
\fancyhead[R]{\small Research draft}
\fancyfoot[C]{\thepage}
\renewcommand{\headrulewidth}{0.3pt}
\setlength{\parskip}{0.3em}
\setlength{\parindent}{1.2em}
\setlist{itemsep=0.2em,topsep=0.4em}
\emergencystretch=2em
\allowdisplaybreaks[2]
\numberwithin{equation}{section}
\newtheorem{theorem}{Theorem}[section]
\newtheorem{lemma}[theorem]{Lemma}
\newtheorem{proposition}[theorem]{Proposition}
\newtheorem{corollary}[theorem]{Corollary}
\theoremstyle{definition}
\newtheorem{definition}[theorem]{Definition}
\newtheorem{example}[theorem]{Example}
\theoremstyle{remark}
\newtheorem{remark}[theorem]{Remark}
\newcommand{\N}{\mathbb N}
\newcommand{\Z}{\mathbb Z}
\newcommand{\R}{\mathbb R}
\newcommand{\ee}{\mathbf e}
\newcommand{\zero}{\mathbf 0}
\newcommand{\one}{\mathbf 1}
\newcommand{\B}{\mathcal B}
\newcommand{\A}{\mathcal A}
\newcommand{\HH}{\mathcal H}
\newcommand{\supp}{\operatorname{supp}}
\newcommand{\Max}{\operatorname{Max}}
\newcommand{\des}{\operatorname{des}}
\newcommand{\lk}{\operatorname{lk}}
\newcommand{\cont}{\operatorname{cont}}
\newcommand{\rk}{\operatorname{rk}}
\newcommand{\lex}{\mathrm{lex}}
\newcommand{\pow}[1]{\langle #1\rangle}
\newcommand{\ind}{\mathbf 1}
\newcommand{\ceil}[1]{\left\lceil #1\right\rceil}
\hypersetup{pdftitle={An explicit shelling for preorder lattice-point posets},
 pdfauthor={AI-assisted research draft},
 pdfsubject={A proposed affirmative solution to Athanasiadis--Chapoton Question 4.6}}

\begin{document}
\begin{center}
{\small\sffamily\color{Accent}ORDER THEORY \; / \; TOPOLOGICAL COMBINATORICS}\\[0.8em]
{\LARGE\bfseries\color{Ink}An explicit shelling for\\[0.15em]
preorder lattice-point posets}\\[0.8em]
{\large A base-exchange solution to Question 4.6}\\[0.9em]
{\normalsize AI-assisted research draft}\\[0.3em]
{\normalsize 20 September 2026}
\end{center}
\vspace{0.3em}
\begin{abstract}
For a finite preorder $\tau$ on $[n]$, let $P_\tau$ be the nonnegative integer
vectors $a$ satisfying $\sum_{i\in I}a_i\le |I|$ for every order ideal $I$ of
$\tau$, ordered coordinatewise. We give an explicit shelling of
$\Delta(P_\tau\setminus\{\zero\})$, and hence of $\Delta(P_\tau)$. This gives an
affirmative answer to both parts of Question 4.6 of Athanasiadis and Chapoton.
The construction groups saturated chains by their maximal vector. Maximal
vectors are ordered decreasingly lexicographically; within each group, words
are ordered lexicographically using a basis-dependent alphabet. The necessary
change of alphabet is controlled by the coordinates from which a unit can be
transferred to an earlier coordinate.

The proof applies more generally to the integer-point poset of every integral
polymatroid. We prove the required exchange property both by tight-set
arguments and, for preorders, by a direct supply-assignment argument. We also
identify every shelling restriction face, obtain ordinary and flag
$h$-enumerators, prove shellability of all rank selections, and show that the
preorder complexes with the minimum removed are contractible but need not be
balls. Exact, definition-level tests cover all 389 labeled preorders on one
through four elements, 126 exchange families in three coordinates, and 58
weighted-preorder cases. The finite computations support the proof; they are
not used to justify any all-size assertion.
\end{abstract}

\smallskip
\noindent\fcolorbox{Ink!35}{Ink!3}{\begin{minipage}{0.95\linewidth}
\small\textbf{Status.} This is a complete proposed mathematical argument, not
an independently refereed or proof-assistant-verified result. The target is
explicitly open in the retrieved primary source. Targeted literature searches
found no resolution of that question; they do not establish priority for every
general lemma proved here. Standard polymatroid and shelling theory is
credited rather than presented as newly invented.
\end{minipage}}

\clearpage
\begingroup
\small
\setlength{\parskip}{0pt}
\tableofcontents
\endgroup
\clearpage

\section{The question and the contribution}
\label{sec:introduction}

Question 4.6 of Athanasiadis and Chapoton \cite[p.~13]{AC} asks:
\begin{quote}
``Is the poset $P_\tau$ shellable for every preorder $\tau$? Is it
Cohen--Macaulay for every preorder $\tau$?''
\end{quote}
Their preceding argument establishes the case in which $\tau$ has a maximum
vertex. The construction below does not require that hypothesis.

\subsection{Why this is a different problem from the manifest entries}

The supplied manifest contains two reports in the same family:
\emph{Reciprocity and Matrix Duality for Preorder Polytopes} and
\emph{A Reflexive Root-Polytope Model for Preorder $h$-Polynomials}. Their
catalogued targets concern reciprocity, matrix duality, and a polyhedral model
for a support enumerator. Neither catalogue entry claims a shelling of the
lattice-point poset or a solution to Question 4.6. This report takes the same
objects in a different direction: the topology of their chain complexes.

The manifest is used as a selection and nonduplication record, not as a source
of proved intermediate theorems. In particular, no assertion from either
listed report is assumed below. The input manifest itself cautions that its
summaries record unrefereed claims rather than audited proofs. The package's
\texttt{manifest\_selection.md} records this boundary explicitly.

\subsection{The main statement}

We use $\N=\{0,1,2,\ldots\}$. For a preorder $\tau$ on $E=[n]$, define
\begin{equation}
 P_\tau=\left\{a\in\N^n:
 a(I):=\sum_{i\in I}a_i\le |I|
 \text{ for every order ideal }I\text{ of }\tau\right\}.
 \label{eq:preorder-poset}
\end{equation}
The order on $P_\tau$ is coordinatewise comparison. A preorder is reflexive
and transitive; antisymmetry is not assumed. An order ideal is a subset closed
under taking smaller elements. Thus equivalent elements of the preorder
always occur together in an ideal, but remain distinct coordinates in
\eqref{eq:preorder-poset}.

\begin{theorem}[Preorder shellability]
\label{thm:main}
For every finite preorder $\tau$, both $\Delta(P_\tau)$ and
$\Delta(P_\tau\setminus\{\zero\})$ are pure shellable complexes. Consequently
they are Cohen--Macaulay over $\Z$ and over every field. If $n\ge1$, the second
complex has dimension $n-1$ and is contractible. Every subposet obtained by
selecting any set of ranks of $P_\tau\setminus\{\zero\}$ is also shellable.
\end{theorem}

The shelling, not just its existence, is part of the theorem. Given the list
of maximal vectors, it is specified by unit transfers, two finite sorting
rules, and multiset words. Its restrictions are explicitly computable.

\subsection{Where the actual work lies}

Every principal interval $[\zero,b]$ is a product of chains and is easy to
shell. That observation alone says nothing about whether shellings of
different principal intervals fit together. A naive lexicographic order of
all chain words can fail even for a rank-two matroid; Section~\ref{sec:examples}
gives a four-word counterexample.

There are two steps that make the construction work. First, exchange implies
that the part of $[\zero,b]$ already covered by earlier maximal vectors is a
union of codimension-one coordinate boxes $[\zero,b-\ee_j]$. Second, these
particular coordinates are put \emph{last} in the alphabet used to sort the
words ending at $b$. Any overlap not already accounted for by an ordinary
word descent then becomes a terminal ridge. Both steps are necessary in the
proof, and both are checked separately by the verifier.

\section{Chains, facets, and restriction faces}
\label{sec:preliminaries}

For $a,b\in\N^s$, write $a\le b$ when $a_i\le b_i$ for every $i$, and set
$|a|=\sum_i a_i$. A finite nonempty set $D\subseteq\N^s$ is a
\emph{downset} if $a\in D$ and $\zero\le u\le a$ imply $u\in D$. Its maximal
elements will be called \emph{bases}. Suppose all bases have the same total
sum $r$. Then every maximal saturated chain has the form
\begin{equation}
 \zero=a^{(0)}<a^{(1)}<\cdots<a^{(r)}=b,
 \qquad a^{(k)}-a^{(k-1)}=\ee_{w_k}.
 \label{eq:chain-word}
\end{equation}
The word $w=w_1\cdots w_r$ has content $b$, meaning that the letter $i$
occurs $b_i$ times. Conversely, every word of content $b$ gives a chain in
$D$, since every prefix vector is at most $b$. Thus there are exactly
\begin{equation}
 N(D)=\sum_{b\in\B}\frac{r!}{\prod_{i=1}^s b_i!}
 \label{eq:facet-count}
\end{equation}
maximal saturated chains, where $\B=\Max(D)$.

\begin{definition}
The order complex $\Delta(Q)$ of a finite poset $Q$ has the elements of $Q$
as vertices and the finite chains as faces, including the empty face. A facet
is an inclusion-maximal face. A complex is \emph{pure of dimension $d$} when
every facet has $d+1$ vertices.
\end{definition}

Our primary complex is
\[
 K(D):=\Delta(D\setminus\{\zero\}).
\]
For a word $w$ as in \eqref{eq:chain-word}, its facet is
$F(w)=\{a^{(1)},\ldots,a^{(r)}\}$. Hence $K(D)$ is pure of dimension $r-1$.
We write $\pow{F}$ for the simplex consisting of all subsets of $F$, and
$\pow{F_1,\ldots,F_m}$ for the complex generated by the indicated facets.

\begin{definition}[Pure shelling]
An order $F_1,\ldots,F_N$ of the facets of a pure $(r-1)$-dimensional complex
is a shelling when, for every $k>1$,
\begin{equation}
 \pow{F_k}\cap\pow{F_1,\ldots,F_{k-1}}
 \label{eq:shell-def}
\end{equation}
is generated by a nonempty collection of the ridges $F_k\setminus\{v\}$,
$v\in F_k$. For $r=1$, a ridge is the empty face. The single empty-facet
complex, of dimension $-1$, is declared shellable.
\end{definition}

If the intersection in \eqref{eq:shell-def} is the union of the ridges
$F_k\setminus\{v\}$ for $v\in T_k$, then a face $G\subseteq F_k$ is new
exactly when $T_k\subseteq G$. The face $T_k$ is the \emph{restriction face}.
For the first facet its restriction face is empty. This description will let
us prove shellability and the enumeration formula simultaneously.

One useful equivalent criterion is the following: for every $i<k$, there is
$j<k$ such that
\begin{equation}
 F_i\cap F_k\ \subseteq\ F_j\cap F_k
       =F_k\setminus\{v\}
 \quad\text{for some }v\in F_k.
 \label{eq:pairwise-criterion}
\end{equation}
Indeed, every face of \eqref{eq:shell-def} belongs to one of the intersections
$F_i\cap F_k$. Criterion~\eqref{eq:pairwise-criterion} says exactly that each
is contained in an old ridge. We use the definition for the main proof and
this criterion for rank selection.

\section{An exchange condition and the overlap lemma}
\label{sec:exchange}

Fix a nonempty finite family $\B\subseteq\N^s$ whose members all have sum $r$,
and let $D=\bigcup_{b\in\B}[\zero,b]$. For the moment, assume the following
specified-insertion exchange condition:
\begin{equation}
 \begin{split}
 &b,c\in\B,\quad c_i>b_i\\[-0.2em]
 &\hspace{1.5em}\Longrightarrow\quad
 \text{there is }j\text{ with }b_j>c_j
 \text{ and }b+\ee_i-\ee_j\in\B.
 \end{split}
 \tag{E}\label{eq:exchange}
\end{equation}
This is not an assumption that \emph{every} surplus coordinate can be used,
and it is not a claim of a simultaneous exchange bijection. Only the stated
existence is needed. Sections~\ref{sec:assignments} and \ref{sec:polymatroids}
prove it for the families to which we apply the theorem.

Use the ordinary order $1<2<\cdots<s$ to compare vectors lexicographically.
A vector $c$ is earlier than $b$ if $c>_{\lex}b$, so bases are listed in
\emph{decreasing} lexicographic order. Define
\begin{equation}
 R(b)=\{j:b_j>0\text{ and }b+\ee_i-\ee_j\in\B
                        \text{ for some }i<j\}.
 \label{eq:active}
\end{equation}
We call these the active coordinates of $b$. The word ``active'' is local to
this construction; no identification with a different matroid activity
statistic is presumed.

\begin{lemma}[Exact overlap]
\label{lem:overlap}
Under \eqref{eq:exchange}, for every $b\in\B$,
\begin{equation}
 [\zero,b]\cap\bigcup_{c>_{\lex}b}[\zero,c]
       =\bigcup_{j\in R(b)}[\zero,b-\ee_j].
 \label{eq:overlap}
\end{equation}
Moreover, $R(b)$ is empty exactly for the first base.
\end{lemma}

\begin{proof}
If $j\in R(b)$, choose $i<j$ with $d=b+\ee_i-\ee_j\in\B$. Then
$d>_{\lex}b$ and $b-\ee_j\le d$, proving the inclusion from right to left.

Conversely, take $a\le b,c$ with $c>_{\lex}b$, and let $i$ be the first
coordinate at which $b$ and $c$ differ. Then $c_i>b_i$, and
\eqref{eq:exchange} gives $j$ with $b_j>c_j$ and
$b+\ee_i-\ee_j\in\B$. Coordinates before $i$ agree, and $j\ne i$, so $j>i$.
Thus $j\in R(b)$. Since the coordinates are integers,
$a_j\le c_j\le b_j-1$, while $a_k\le b_k$ for all $k\ne j$.
Consequently $a\le b-\ee_j$.

The first base has no active coordinate, because any witnessing transfer
would produce an earlier base. For every later base, applying the preceding
argument to any earlier base supplies an active coordinate.
\end{proof}

\begin{remark}[The weaker hypothesis actually used]
\label{rem:weaker}
For the shelling theorem, it is enough to be given an ordering of equidegree
bases for which each noninitial overlap has the form on the right of
\eqref{eq:overlap}, with a nonempty set of coordinates. Neither a global
lexicographic order nor the full exchange condition is then needed. Exchange
provides a concrete way to obtain this overlap property and to compute its
coordinate set. We state the proof in the explicit exchange setting to avoid
hiding this step in a multicomplex-shellability assertion.
\end{remark}

\section{The base-block shelling}
\label{sec:shelling}

For each $b$, choose the following total order $\prec_b$ on the letters:
\begin{equation}
 [s]\setminus R(b)\quad\text{first, followed by}\quad R(b),
 \label{eq:alphabet}
\end{equation}
with the usual order within each part. Zero-content letters may be retained
in the alphabet; they do not occur in any word of content $b$.

The complete facet order is now specified. Visit bases in decreasing ordinary
lexicographic order. For each base $b$, visit all words of content $b$ in
\emph{increasing} lexicographic order with respect to $\prec_b$. Replace each
word by its prefix-vector facet $F(w)$. This gives every facet exactly once:
the terminal vector identifies its block, and the successive differences
identify its word.

Write
\begin{equation}
 \operatorname{Des}_b(w)=
 \{k\in\{1,\ldots,r-1\}:w_{k+1}\prec_b w_k\}.
 \label{eq:descent}
\end{equation}
Here the comparisons inside words use $\prec_b$, not necessarily the natural
order on the indices. Bases themselves are still ordered using the natural
coordinate order.

\subsection{What earlier words in one block cover}

\begin{lemma}[Descent ridges within a box]
\label{lem:word}
Fix a base $b$ and any total alphabet order $\prec$. Among the words of
content $b$, put the facets in increasing lexicographic order. For the facet
$F(w)$, its intersection with the earlier facets in this same block is
precisely the union of the ridges
\begin{equation}
 F(w)\setminus\{a^{(k)}\},\qquad
             w_{k+1}\prec w_k.
 \label{eq:word-ridges}
\end{equation}
If there are no descents, there are no earlier facets in the block.
\end{lemma}

\begin{proof}
Consider a face $G\subseteq F(w)$. The retained prefix vectors determine the
content of each segment between successive retained ranks. Add ranks $0$ and
$r$ as segment boundaries, whether or not the terminal vertex belongs to $G$.
The endpoints $\zero$ and $b$ are fixed for every chain in the block.

A word of content $b$ contains $G$ exactly when each of these segments has its
prescribed content. The lexicographically first such word is obtained by
sorting every segment weakly increasingly. Therefore there is an earlier word
containing $G$ exactly when at least one segment of $w$ is not weakly
increasing. Such a segment contains an adjacent descent; its rank $k$ is not
retained in $G$. Conversely, if $G$ omits a descent rank $k$, swapping the
adjacent descending letters gives an earlier word, changes only the rank-$k$
prefix vector, and keeps every vertex of $G$.

Thus $G$ is old within the block exactly when it misses at least one descent
vertex, which is precisely the union in \eqref{eq:word-ridges}. The only
word with no descents is the sorted word, the first word of the block.
\end{proof}

\subsection{What all earlier blocks cover}

\begin{lemma}[The last active letter]
\label{lem:last-active}
Let $b$ be a noninitial base, let $w$ have content $b$, and set
\begin{equation}
 t(w)=\max\{k:w_k\in R(b)\}.
 \label{eq:last-active}
\end{equation}
The faces of $F(w)$ appearing in earlier base blocks are exactly the subsets
of $\{a^{(1)},\ldots,a^{(t(w)-1)}\}$.
\end{lemma}

\begin{proof}
By Lemma~\ref{lem:overlap}, a prefix vector $a^{(k)}$ is in an earlier
principal box if and only if some active coordinate is not yet exhausted,
that is, $a^{(k)}_j<b_j$ for some $j\in R(b)$. This is equivalent to
$k<t(w)$.

For a nonempty face $G$ of $F(w)$, let $a^{(k)}$ be its largest vertex. If
$k<t(w)$, that vertex lies below some earlier base $c$. Every other vertex
of $G$ lies below it, so the whole face lies in $[\zero,c]$. A chain in this
box extends to a saturated chain ending at $c$, whose block is already
complete. Thus $G$ is old. If $k\ge t(w)$, its largest vertex belongs to no
earlier box, so $G$ is not old in an earlier block. The empty face is old
because there is an earlier block. This proves the stated equality.
\end{proof}

\subsection{The intersection identity and shellability}

\begin{theorem}[Explicit shelling and restriction face]
\label{thm:abstract}
Let $D=\bigcup_{b\in\B}[\zero,b]$, where $\B$ is a nonempty equidegree family
satisfying \eqref{eq:exchange}. The base-block order is a shelling of $K(D)$.
For its facet $F(w)$ ending at $b$, the restriction face is
\begin{equation}
 \mathcal R(w)=
 \{a^{(k)}:k\in\operatorname{Des}_b(w)\}
 \ \cup\
 \begin{cases}
   \{b\},&w_r\in R(b),\\
   \varnothing,&w_r\notin R(b).
 \end{cases}
 \label{eq:restriction}
\end{equation}
For $r=0$ there is one empty facet, with empty restriction face.
\end{theorem}

\begin{proof}
Assume $r\ge1$. Within the first block, Lemma~\ref{lem:word} gives exactly the
descent ridges. There are no earlier blocks and $R(b)=\varnothing$, so
\eqref{eq:restriction} follows there.

Now consider a noninitial base $b$. Earlier facets split into those in its
own block and those in earlier blocks. The former give exactly the descent
ridges. The latter, by Lemma~\ref{lem:last-active}, give the simplex on the
prefix vertices of ranks below $t=t(w)$.

If $t<r$, then $w_t\in R(b)$ and $w_{t+1}\notin R(b)$. By
\eqref{eq:alphabet}, this is a descent at rank $t$. Every face from an earlier
block misses $a^{(t)}$ and is therefore already contained in the descent ridge
$F(w)\setminus\{a^{(t)}\}$. There is no extra contribution from earlier blocks.

If $t=r$, then the faces from earlier blocks form exactly the terminal ridge
$F(w)\setminus\{b\}$. This case is equivalent to $w_r\in R(b)$. Combining the
two cases, the full old intersection is
\begin{equation}
 \bigcup_{k\in\operatorname{Des}_b(w)}
       \pow{F(w)\setminus\{a^{(k)}\}}
 \quad\cup\quad
 \begin{cases}
  \pow{F(w)\setminus\{b\}},&w_r\in R(b),\\
  \text{no additional faces},&w_r\notin R(b).
 \end{cases}
 \label{eq:full-intersection}
\end{equation}
This is a union of ridges, not merely a subcomplex of such a union.

Except for the very first facet, the union contains at least one ridge. A
nonfirst word in a block has a descent. For the first word of a noninitial
block, the sorted word ends in an active letter, since $R(b)\ne\varnothing$
and every active coordinate has positive content. Its terminal ridge is old.
Thus \eqref{eq:full-intersection} proves shellability. A face is absent from
that union exactly when it contains every omitted vertex indexing its ridges;
these vertices are exactly \eqref{eq:restriction}.
\end{proof}

\begin{remark}[A directly checkable witness]
When $w_r=j\in R(b)$, choose $i<j$ with $d=b+\ee_i-\ee_j\in\B$.
Replacing the final letter of $w$ by $i$ produces a word ending at $d$ in an
earlier block; its facet shares exactly the terminal ridge. For an internal
descent, swapping the two letters gives the corresponding earlier facet in
the same block. Thus every ridge in \eqref{eq:full-intersection} has an
explicit earlier-facet witness.
\end{remark}

\section{Why preorder posets satisfy exchange}
\label{sec:assignments}

We give a combinatorial proof directly from the defining inequalities. It
makes no assumption about a unique maximal equivalence class of the preorder.
For $A\subseteq E$, write
\[
 \downarrow A=\{u\in E:u\le_\tau a\text{ for some }a\in A\}.
\]
This is an ideal, and $A\subseteq\downarrow A$.

\begin{lemma}[Equivalent constraints]
\label{lem:coverage-constraints}
A vector $a\in\N^n$ belongs to $P_\tau$ exactly when
\begin{equation}
 a(A)\le |\downarrow A|\qquad(A\subseteq E).
 \label{eq:coverage}
\end{equation}
\end{lemma}

\begin{proof}
The ideal inequalities imply
$a(A)\le a(\downarrow A)\le|\downarrow A|$ by nonnegativity. Conversely,
$\downarrow I=I$ for an ideal $I$, so \eqref{eq:coverage} includes every
defining inequality.
\end{proof}

\subsection{A supply interpretation and purity}

Place one labeled supply token at every $u\in E$. A token at $u$ may be
assigned to coordinate $i$ precisely when $u\le_\tau i$. Tokens may also be
left unused. The occupancy vector records the number assigned to each
coordinate.

\begin{lemma}[Supply representation]
\label{lem:supply}
The possible occupancy vectors are exactly $P_\tau$. Its maximal vectors are
exactly the occupancy vectors of complete assignments, and all have total
sum $n$.
\end{lemma}

\begin{proof}
For a proposed occupancy vector $a$, create $a_i$ demand copies at each
coordinate $i$. Join each demand copy at $i$ to all supply tokens in
$\downarrow\{i\}$. A matching saturating these demand copies is precisely an
assignment realizing $a$.

For completeness, the finite matching criterion is that every set of demand
vertices have at least as many supply neighbors. Necessity is immediate.
For sufficiency, take a maximum matching. If a demand vertex is unmatched,
follow alternating paths from it. If an unmatched supply vertex is reached,
the matching can be enlarged. Otherwise, all reached supply vertices are
matched to reached demand vertices, while the initial demand vertex is
unmatched. The reached demand set then has fewer supply neighbors than its
cardinality, a contradiction.

For any set of demand copies, let $A$ be its set of coordinate labels. Its
cardinality is at most $a(A)$, and its supply neighborhood is $\downarrow A$.
Thus \eqref{eq:coverage} ensures the matching criterion. Conversely, tokens
assigned to coordinates in $A$ are distinct tokens in $\downarrow A$, so
$a(A)\le|\downarrow A|$ is necessary.
Lemma~\ref{lem:coverage-constraints} proves the first assertion.

If fewer than $n$ tokens are used, choose an unused token at $u$ and assign
it to coordinate $u$, which is allowed by reflexivity. This strictly increases
the occupancy vector. Hence a maximal vector has total $n$. Every vector of
total $n$ is maximal because the inequality for $E$ forbids a larger one.
\end{proof}

\subsection{Exchange by rerouting assignments}

\begin{lemma}[Specified insertion]
\label{lem:assignment-exchange}
The family $\B_\tau=\Max(P_\tau)$ satisfies \eqref{eq:exchange}.
\end{lemma}

\begin{proof}
Take complete assignments realizing bases $b,c$. For every supply token,
draw a directed edge from its destination in the $b$-assignment to its
destination in the $c$-assignment, retaining the token as the edge label.
Loops and parallel edges are allowed. At coordinate $v$,
\[
 \operatorname{outdeg}(v)-\operatorname{indeg}(v)=b_v-c_v.
\]
Suppose $c_i>b_i$. Let $S$ be the set of vertices from which a directed path
to $i$ exists. No directed edge enters $S$ from outside: its initial vertex
would then also reach $i$. Therefore
\[
 \sum_{v\in S}(b_v-c_v)
   =\#\{\text{edges leaving }S\}\ge0.
\]
The summand at $i$ is negative, so some $j\in S$ satisfies $b_j>c_j$.
Choose a simple directed path from $j$ to $i$.

For each edge on this path, reassign its labeled token from its
$b$-destination to its $c$-destination. Every new destination is allowed,
since it was used in the $c$-assignment. A simple path uses distinct edges
and hence distinct tokens. At each internal vertex, one token leaves and
one arrives; the occupancy decreases by one at $j$ and increases by one at
$i$. All tokens remain assigned. The resulting base is
$b+\ee_i-\ee_j$, as required.
\end{proof}

Lemmas~\ref{lem:supply} and \ref{lem:assignment-exchange}, followed by
Theorem~\ref{thm:abstract}, prove the shellability assertion of
Theorem~\ref{thm:main}. Notice that the supply representation is an integral
statement: it identifies actual integer vectors and actual matchings, not
only convex hulls.

\section{Extension to integral polymatroids}
\label{sec:polymatroids}

The matching proof is useful for preorders, but its full generality is
unnecessary. A separate tight-set argument shows that the shelling applies
to every integral polymatroid. We include the proof so that no unverified
exchange convention is being imported. The setting is classical; see
Herzog and Hibi \cite{HH}.

Let $f:2^{[s]}\to\N$ be normalized, monotone, and submodular:
\begin{equation}
 \begin{gathered}
 f(\varnothing)=0,\qquad
 A\subseteq B\Longrightarrow f(A)\le f(B),\\
 f(A)+f(B)\ge f(A\cup B)+f(A\cap B).
 \end{gathered}
 \label{eq:rank-function}
\end{equation}
Define its integer independence poset
\begin{equation}
 D_f=\{a\in\N^s:a(A)\le f(A)\text{ for all }A\subseteq[s]\}.
 \label{eq:integral-pm}
\end{equation}
It is a finite downset, since $a_i\le f(\{i\})$. For $a\in D_f$, call a set
$A$ \emph{$a$-tight} if $a(A)=f(A)$.

\begin{lemma}[Tight-set closure]
\label{lem:tight}
The union and intersection of two $a$-tight sets are $a$-tight.
\end{lemma}

\begin{proof}
The vector sum $a(\cdot)$ is modular, so
\[
\begin{split}
 f(A)+f(B)&=a(A)+a(B)\\
 &=a(A\cup B)+a(A\cap B)\\
 &\le f(A\cup B)+f(A\cap B)\le f(A)+f(B).
\end{split}
\]
All inequalities are equalities, proving the assertion.
\end{proof}

\begin{lemma}[Purity and exchange from tight sets]
\label{lem:pm-exchange}
Every maximal vector of $D_f$ has total $r=f([s])$, and its base family
satisfies \eqref{eq:exchange}.
\end{lemma}

\begin{proof}
If a feasible vector $a$ is maximal, $a+\ee_i$ is infeasible for each $i$.
Because all slacks are integers, there is an $a$-tight set $A_i$ containing
$i$. Their union is $[s]$, and Lemma~\ref{lem:tight} shows that it is tight.
Thus $|a|=f([s])$. Conversely, a feasible vector of that total cannot be
increased, by the inequality for $[s]$.

Let $b,c$ be bases with $c_i>b_i$. The collection of $b$-tight sets containing
$i$ is nonempty, since $[s]$ is tight. Let $T$ be its intersection. The set
$T$ is tight and contains $i$. Since $c(T)\le f(T)=b(T)$, some $j\in T$
satisfies $b_j>c_j$.

Set $d=b+\ee_i-\ee_j$. It is nonnegative because $b_j\ge1$. A constraint
$d(A)\le f(A)$ could fail only if $i\in A$, $j\notin A$, and $A$ was
$b$-tight: every nontight constraint has integral slack at least one. But
by the definition of $T$, every tight set containing $i$ also contains $j$.
No constraint fails. Finally, $|d|=r$, so $d$ is a base.
\end{proof}

\begin{corollary}[Integral-polymatroid shellability]
\label{cor:all-pm}
For every $f$ satisfying \eqref{eq:rank-function}, the complexes
$\Delta(D_f\setminus\{\zero\})$ and $\Delta(D_f)$ are pure shellable and
Cohen--Macaulay over $\Z$ and every field. The base-block order and
restriction formula \eqref{eq:restriction} apply without change.
\end{corollary}

\begin{proof}
Apply Lemma~\ref{lem:pm-exchange} and Theorem~\ref{thm:abstract}; the
Cohen--Macaulay consequence is proved in Section~\ref{sec:topology}.
\end{proof}

For a preorder, $f(A)=|\downarrow A|$ is a coverage function. Indeed,
$\downarrow(A\cup B)=\downarrow A\cup\downarrow B$ and
$\downarrow(A\cap B)\subseteq\downarrow A\cap\downarrow B$, which prove
submodularity by cardinalities. This supplies a second route from
\eqref{eq:preorder-poset} to the shelling.

\subsection{Weighted preorders and arbitrary supply neighborhoods}

Given capacities $c_u\in\N$, define
\begin{equation}
 P_{\tau,c}=\{a\in\N^n:a(I)\le c(I)
                         \text{ for every ideal }I\}.
 \label{eq:weighted}
\end{equation}
The rank function is $f_c(A)=c(\downarrow A)$, which is again normalized,
monotone, integral, and submodular. Its total rank is $\sum_u c_u$. Thus the
same shelling works, including when some capacities vanish. Equivalently,
replace the supply token at $u$ by $c_u$ distinct copies in the matching
proof. The contractibility assertion for the unweighted case is not being
silently extended to all zero-capacity cases.

More generally, let $U$ be a finite supply set and let each token $u$ have
a nonempty allowed destination set $N(u)\subseteq[s]$. All partial
assignments have occupancies in the coverage polymatroid with rank
\[
 f(A)=|\{u:N(u)\cap A\ne\varnothing\}|.
\]
The matching proof identifies its integer points with those occupancies,
and unused tokens can be assigned to any allowed destination. The directed
rerouting proof establishes exchange verbatim. Thus neither transitivity nor
reflexivity is essential for this larger supply model; they were what
identified the original ideal inequalities with the preorder supply model.

\section{Topology and rank selection}
\label{sec:topology}

\subsection{Why shellability gives the required Cohen--Macaulay property}

We use the homological definition: a pure finite complex $K$ is
Cohen--Macaulay over $\Z$ if, for every face $\sigma$ (including the empty
face),
\begin{equation}
 \widetilde H_i(\lk_K\sigma;\Z)=0
       \quad\text{for all }i<\dim(\lk_K\sigma).
 \label{eq:cm}
\end{equation}
Here the link consists of faces disjoint from $\sigma$ whose union with
$\sigma$ is a face. Over a field, the same vanishing condition is equivalent
to the usual Cohen--Macaulay condition on the Stanley--Reisner ring. The
standard shelling implications and the order-homotopy tools used below are
reviewed by Wachs \cite{Wachs}. We spell out the relevant implications.

\begin{lemma}
\label{lem:shell-cm}
Every pure shellable finite complex satisfies \eqref{eq:cm}, has free abelian
homology, and is Cohen--Macaulay over every field.
\end{lemma}

\begin{proof}
The facets containing a fixed face $\sigma$, with $\sigma$ deleted, give a
shelling of the link in their inherited order. To see this, apply
\eqref{eq:pairwise-criterion} to two facets containing $\sigma$. Its witness
intersection contains $\sigma$, so the witness facet also contains $\sigma$,
and its omitted vertex is outside $\sigma$. The same criterion therefore
holds after deleting $\sigma$.

It remains to examine an arbitrary pure shellable $d$-complex $L$. It starts
with a simplex. Each later simplex is attached along a nonempty collection
of boundary facets. A proper such collection is contractible: choose a
boundary facet not in the collection; the vertex opposite it belongs to all
facets in the collection, so their union is a cone. The full collection is
the boundary sphere. The reduced Mayer--Vietoris sequence shows inductively
that no homology below dimension $d$ is created. In the full-boundary case,
the top homology acquires one free summand; in the proper-boundary case it is
unchanged. For $d=0$, the complex is a collection of vertices and the same
conclusion follows directly. Dimension $-1$ uses the empty-face convention.
Applying this to every link proves \eqref{eq:cm}. Free homology and the
universal coefficient theorem give the identical lower-dimensional
vanishing over every field.
\end{proof}

Since $\zero$ is a least element, $\Delta(D)$ is the cone on $K(D)$.
Appending that common vertex to each facet preserves its shelling order.
Its Cohen--Macaulay property then follows from Lemma~\ref{lem:shell-cm}.
This is stronger than observing that the cone is contractible: a
contractible complex need not be Cohen--Macaulay.

\subsection{Every rank selection is shellable}

For $S\subseteq[r]$, let $D_S=\{a\in D:|a|\in S\}$. A facet of
$\Delta(D_S)$ is obtained by retaining the ranks in $S$ from a full facet.
List the distinct retained facets in order of their first appearance in the
base-block shelling.

\begin{proposition}
\label{prop:rankselection}
This first-appearance order is a shelling of $\Delta(D_S)$.
\end{proposition}

\begin{proof}
All facets have $|S|$ vertices: every chain can be extended to a full
saturated chain in $D$. Let $G$ be a newly appearing retained facet, and
choose the earliest full facet $F$ yielding it. For an earlier retained facet
$G'$, choose an earlier full facet $F'$ yielding it. The full shelling
criterion provides an earlier full facet $H$ with
\[
 F'\cap F\subseteq H\cap F=F\setminus\{v\}.
\]
If the rank of $v$ were outside $S$, the restriction of $H$ to $S$ would
already equal $G$, contradicting the choice of $F$. Thus $v$ has rank in $S$.
The restriction of $H$ is an earlier retained facet and meets $G$ in exactly
$G\setminus\{v\}$, containing $G'\cap G$. The criterion applies.
\end{proof}

\subsection{Removing the minimum does not destroy contractibility here}

\begin{proposition}
\label{prop:contractible}
For $n\ge1$, $\Delta(P_\tau\setminus\{\zero\})$ is contractible.
\end{proposition}

\begin{proof}
Every binary vector belongs to $P_\tau$, because its sum on an ideal $I$ is
at most $|I|$. On the nonzero part define
\[
 q(a)_i=\begin{cases}1,&a_i>0,\\0,&a_i=0.\end{cases}
\]
This is an order-preserving map satisfying $q(a)\le a$ and $q^2=q$. Its image
is the poset of nonzero binary vectors, which has maximum $\one$.

The order maps $q$ and the identity induce homotopic maps on order complexes.
One can see this directly by the order-preserving map
$H:(P_\tau\setminus\{\zero\})\times\{0<1\}\to
P_\tau\setminus\{\zero\}$ with $H(a,0)=q(a)$ and $H(a,1)=a$; the standard
prism triangulation on each chain gives the homotopy. The map $q$ fixes its
image, so the image inclusion and $q$ are homotopy inverses. The image order
complex is a cone at $\one$, hence contractible.
\end{proof}

This proof uses no shellability. It independently checks the vanishing of the
top reduced homology in the special preorder family. Together with the
preceding results it completes all assertions of Theorem~\ref{thm:main}.
For the empty preorder, $P_\tau=\{\zero\}$; its order complex is a point and
its nonzero part has the single empty face. Contractibility of the latter is
not asserted.

\section{Ordinary and flag \texorpdfstring{$h$}{h}-enumerators}
\label{sec:enumeration}

For a pure $(r-1)$-complex $K$, define
\begin{equation}
 h(K;z)=\sum_{G\in K}z^{|G|}(1-z)^{r-|G|}.
 \label{eq:hdef}
\end{equation}
If $f_{j-1}$ is the number of faces of cardinality $j$, this is the standard
face-to-$h$ transformation with $f_{-1}=1$. The shelling partitions the face
set into Boolean intervals $[\mathcal R(w),F(w)]$. Such an interval
contributes exactly $z^{|\mathcal R(w)|}$ to \eqref{eq:hdef}, by the binomial
theorem.

\begin{corollary}[A nonnegative descent formula]
\label{cor:hformula}
For $D$ satisfying the hypotheses of Theorem~\ref{thm:abstract},
\begin{equation}
 \boxed{\quad h(K(D);z)=
 \sum_{b\in\B}\ \sum_{\cont(w)=b}
 z^{\des_{\prec_b}(w)+\ind_{\{w_r\in R(b)\}}}.\quad}
 \label{eq:hformula}
\end{equation}
When $r=0$, the unique empty word is assigned exponent zero, so the sum is
$1$. The same polynomial is $h(\Delta(D);z)$, because coning preserves the
$h$-polynomial. In particular its coefficients are nonnegative integers,
its constant coefficient is one, and its value at $z=1$ is \eqref{eq:facet-count}.
\end{corollary}

The formula counts facets, not lattice points. The extra terminal indicator
is essential: counting only ordinary descents within the blocks would assign
a constant term to the first word of \emph{every} block, rather than only to
the very first facet.

\subsection{The flag refinement}

Color a nonzero vector $a$ by its rank $|a|\in[r]$. Every facet contains one
vertex of each color. For $T\subseteq[r]$, let $f_T$ be the number of faces
whose rank set is exactly $T$, and define
\[
 h_S=\sum_{T\subseteq S}(-1)^{|S|-|T|}f_T.
\]
For a word $w$ ending at $b$, put
\begin{equation}
 A(w)=\operatorname{Des}_b(w)\cup
       \begin{cases}\{r\},&w_r\in R(b),\\
                     \varnothing,&w_r\notin R(b).
       \end{cases}
 \label{eq:restriction-ranks}
\end{equation}
For the empty word ($r=0$), set $A(w)=\varnothing$.

\begin{proposition}[Flag $h$ by restriction ranks]
\label{prop:flagh}
For every $S\subseteq[r]$,
\begin{equation}
 h_S=\#\{w:\cont(w)\in\B,\ A(w)=S\}.
 \label{eq:flagh}
\end{equation}
Consequently, for any rank set $S$,
\begin{equation}
 h(\Delta(D_S);z)=\sum_{T\subseteq S}h_Tz^{|T|}.
 \label{eq:selected-h}
\end{equation}
\end{proposition}

\begin{proof}
The interval $[\mathcal R(w),F(w)]$ contains one face of rank set $T$ if
$A(w)\subseteq T$, and none otherwise. Hence
$f_T=\sum_{U\subseteq T}\#\{w:A(w)=U\}$. Boolean-lattice inversion gives
\eqref{eq:flagh}. Rank selection leaves all $f_T$, and hence all $h_T$, with
$T\subseteq S$ unchanged. Grouping them by cardinality gives
\eqref{eq:selected-h}.
\end{proof}

\subsection{Which polynomial this is, and which it is not}

The support enumerator
\[
 h_\tau^{\mathrm{supp}}(z)=\sum_{a\in P_\tau}z^{|\supp(a)|}
\]
is the preorder $h$-polynomial studied in the manifest's root-polytope report.
It is generally different from \eqref{eq:hformula}. Support-enumerator
duality is treated by Dai, Hou, Liu, Thawinrak, and Wang \cite{DHLTW}; it is
not an input to the shelling proof.

For clarity, $\mathcal Q_\tau$ denotes the real polytope obtained by
replacing $\N^n$ in \eqref{eq:preorder-poset} by $\R_{\ge0}^n$.
Its Ehrhart $h^*$-polynomial is the numerator determined by
\[
 \sum_{m\ge0}\#(m\mathcal Q_\tau\cap\Z^n)z^m
       =\frac{h^*(\mathcal Q_\tau;z)}{(1-z)^{n+1}}.
\]
Here the dimension is $n$, since the polytope contains the unit cube.
Ehrhart--Zeta duality \cite[Theorem~1.5 and the proof of Proposition~4.5]{AC}
identifies $h(\Delta(P_\tau);z)$ with $h^*(\mathcal Q_{\tau^*};z)$, where
$\tau^*$ is the opposite preorder. With that \emph{imported} identification,
\eqref{eq:hformula} is also a descent formula for the latter polynomial.
The shelling proof and the flag enumerator do not depend on this import.
This connection is not a claim to reprove the manifest's enumerative results.

In the unweighted preorder case the coefficient of $z^n$ in
\eqref{eq:hformula} vanishes, consistently with
Proposition~\ref{prop:contractible}. A direct check is also instructive: to
have $n-1$ strict descents in a word of length $n$ on $n$ letters, all letters
must occur once. To obtain the additional terminal contribution, the last
letter must be active. But $1$ is never active by \eqref{eq:active}; a strictly
decreasing word in $\prec_b$ ends in a nonactive letter if it contains every
letter. Hence the exponent cannot be $n$.

\section{Examples, failed shortcuts, and boundary cases}
\label{sec:examples}

\subsection{A preorder without a maximum}

Let $\tau$ have $1<_\tau2$ and $1<_\tau3$, with $2$ and $3$ incomparable.
Its ideals give
\[
 a_1\le1,\qquad a_1+a_2\le2,\qquad a_1+a_3\le2,
 \qquad a_1+a_2+a_3\le3.
\]
Its bases, in the chosen order, are
\[
 b^{(1)}=(1,1,1),\qquad b^{(2)}=(0,2,1),\qquad b^{(3)}=(0,1,2).
\]
Their active sets are $\varnothing$, $\{2\}$, and $\{3\}$, respectively.
For the middle base the alphabet changes to $1\prec3\prec2$. This is a small
instance in which the basis-dependent ordering is visible.

\begin{table}[htbp]
\centering\small
\begin{tabular}{@{}rccccc@{}}
\toprule
Step & Base & $R(b)$ & Alphabet & Word & Restriction ranks \\
\midrule
1 & $111$ & $\varnothing$ & $1\prec 2\prec 3$ & $123$ & $\varnothing$ \\
2 & $111$ & $\varnothing$ & $1\prec 2\prec 3$ & $132$ & $\{2\}$ \\
3 & $111$ & $\varnothing$ & $1\prec 2\prec 3$ & $213$ & $\{1\}$ \\
4 & $111$ & $\varnothing$ & $1\prec 2\prec 3$ & $231$ & $\{2\}$ \\
5 & $111$ & $\varnothing$ & $1\prec 2\prec 3$ & $312$ & $\{1\}$ \\
6 & $111$ & $\varnothing$ & $1\prec 2\prec 3$ & $321$ & $\{1,2\}$ \\
7 & $021$ & $\{2\}$ & $1\prec 3\prec 2$ & $322$ & $\{3\}$ \\
8 & $021$ & $\{2\}$ & $1\prec 3\prec 2$ & $232$ & $\{1,3\}$ \\
9 & $021$ & $\{2\}$ & $1\prec 3\prec 2$ & $223$ & $\{2\}$ \\
10 & $012$ & $\{3\}$ & $1\prec 2\prec 3$ & $233$ & $\{3\}$ \\
11 & $012$ & $\{3\}$ & $1\prec 2\prec 3$ & $323$ & $\{1,3\}$ \\
12 & $012$ & $\{3\}$ & $1\prec 2\prec 3$ & $332$ & $\{2\}$ \\
\bottomrule
\end{tabular}

\caption{The complete shelling for the three-element preorder with
$1<2$ and $1<3$. Base strings denote coordinate vectors.}
\label{tab:worked}
\end{table}

There are 12 lattice points including $\zero$. The order complex of its
nonzero part has face counts
\[
 (f_{-1},f_0,f_1,f_2)=(1,11,22,12),
\]
so direct substitution in \eqref{eq:hdef} gives
\[
 (1-z)^3+11z(1-z)^2+22z^2(1-z)+12z^3
       =1+8z+3z^2.
\]
The same answer is obtained by counting the restriction sizes in
Table~\ref{tab:worked}. In contrast, direct support counting gives
$h_\tau^{\mathrm{supp}}(z)=1+5z+5z^2+z^3$. These distinct answers illustrate
why the two $h$-polynomials must not be conflated.

\subsection{Why one global word order is insufficient}

Consider the equidegree family
\[
 \B=\{(1,0,1),(0,1,1)\}.
\]
It satisfies \eqref{eq:exchange} and is the rank-two matroid with element 3
in every base. Global natural lexicographic order on its words is
\[
 13,\quad23,\quad31,\quad32.
\]
The first two facets are the disjoint edges
\[
 \{\ee_1,\ee_1+\ee_3\},\qquad
 \{\ee_2,\ee_2+\ee_3\}.
\]
Their intersection has no vertex, so this is not a shelling of a pure
one-dimensional complex. The base-block construction instead gives
\[
 13,\quad31,\quad32,\quad23,
\]
which builds a path one edge at a time. This example is an illustration of
the general theorem, not itself an unweighted preorder poset of the form
\eqref{eq:preorder-poset}.

\subsection{Why \texorpdfstring{$M$}{M}-shellability cannot be substituted for the proof}

Alizadeh, Goodarzi, and Yassemi \cite{AGY} establish $M$-shellability of
discrete polymatroids. Their notion partitions a monomial order ideal into
intervals that are products of chains, with initial unions forming order
ideals. It is important not to read that statement as automatically asserting
ordinary simplicial shellability of the chain complex.

For example, the pure monomial downset
\[
 \{1,x,x^2,y,y^2\}
\]
has an $M$-shelling consisting of $[1,x^2]$ followed by $[y,y^2]$. After its
minimum is removed, however, the order complex consists of the two disjoint
edges $x<x^2$ and $y<y^2$, so it is not pure shellable. This is not a
counterexample to the theorem of \cite{AGY}: the downset is not a discrete
polymatroid. It shows that the \emph{logical implication from the weaker
notion alone} would be invalid.

Ficarra's work \cite{Ficarra} studies linear quotients and shellability of
simplicial multicomplexes, another closely related setting. Rather than assume
that a terminology match settles the present topological question, we proved
the precise overlap identity and its conversion into an ordinary shelling.
The broader lemmas may have prior formulations in this literature; no
priority claim for them is essential to the solution of the selected question.

\subsection{Contractible does not mean a ball}

Take the universal preorder on three elements: every element is equivalent
to every other. Then
\[
 P_\tau=\{a\in\N^3:|a|\le3\}.
\]
The edge of its nonzero order complex corresponding to
$\ee_1<2\ee_1$ lies in the three triangles with terminal vertices
$3\ee_1$, $2\ee_1+\ee_2$, and $2\ee_1+\ee_3$. A triangulated two-dimensional
ball or sphere cannot have three triangles incident with one edge. Thus the
complex in Theorem~\ref{thm:main} need not be a ball, even though it is
contractible and shellable. No manifold or polyhedral-boundary conclusion is
claimed.

Every order complex is flag: pairwise comparable vertices form a chain. But
this does not yield a flag simplicial \emph{polytope} realizing the support
polynomial. The objects, dimensions, and enumerators are different.

\section{Reproducible exact verification}
\label{sec:verification}

The archive includes a Python 3.10+ verifier using only the standard library.
Its most important tests do not assume the shelling theorem. For each facet
in the proposed order, the verifier enumerates every subset of its actual
prefix-vector vertex set. It maintains the set of all previously appearing
faces, extracts the old ridges, and checks that every old face lies in one
of those ridges. It also checks that a face is new if and only if it contains
the predicted restriction face. Thus both the definition of shellability and
the stronger interval-partition assertion are checked directly.

\subsection{Exhaustive ranges and recorded results}

Every reflexive relation on a labeled $n$-element set is generated by choosing
the $n(n-1)$ off-diagonal bits; transitivity is then tested exactly. This
produces the exhaustive preorder counts in Table~\ref{tab:verification}.
For each preorder, the lattice-point set is independently generated from all
of its ideal inequalities. Its maximal vectors are compared with the
occupancies of all complete labeled-supply assignments.

\begin{table}[htbp]
\centering
\begin{tabular}{@{}rrrr@{}}
\toprule
Ground-set size & Preorders & Facets checked & Faces counted\\
\midrule
1 & 1 & 1 & 2\\
2 & 4 & 12 & 32\\
3 & 29 & 411 & 1\,516\\
4 & 355 & 33\,216 & 172\,566\\
\midrule
Total & 389 & 33\,640 & 174\,116\\
\bottomrule
\end{tabular}
\caption{Exact tests of all labeled preorders through size four. Faces
include the empty face and are counted separately for each preorder.}
\label{tab:verification}
\end{table}

The abstract test examines every nonempty subfamily of all weak compositions
of $r$ into three coordinates, for $r=1,2,3$. It tests the full
specified-insertion condition for every ordered pair of bases and every
eligible coordinate; only the families passing that condition are then
subjected to the shelling tests.

\begin{table}[htbp]
\centering
\begin{tabular}{@{}rrrr@{}}
\toprule
Rank & Base families examined & Exchange families & Facets checked\\
\midrule
1 & 7 & 7 & 12\\
2 & 63 & 29 & 126\\
3 & 1\,023 & 90 & 1\,104\\
\midrule
Total & 1\,093 & 126 & 1\,242\\
\bottomrule
\end{tabular}
\caption{Exhaustive equidegree-family tests in three coordinates.}
\label{tab:abstract-tests}
\end{table}

For each of the 29 preorders on three labels, the two capacity vectors
$(2,1,1)$ and $(0,2,2)$ are also tested. These are 58 weighted cases, with
2\,020 facets checked. The empty preorder is tested separately with the
empty-face convention. No randomness is used in any of these tests.

For every tested complex, its ordinary $h$-polynomial is computed in two
ways: by restriction sizes, and by the direct transform of all distinct
faces. Its complete flag $h$-vector is likewise computed both by restriction
rank sets and by Boolean inversion of the flag face counts. All comparisons
passed. In each unweighted nonempty preorder case, the top ordinary
$h$-coefficient was zero, independently checking the consequence of
contractibility.

\subsection{What the tests do not establish}

Finite testing cannot prove the theorem for arbitrary rank or size, and no
such inference is used. The verifier is not a proof assistant. It does not
formalize the tight-set argument, compute all link homology groups over
$\Z$, establish bibliographic priority, or verify unrelated assertions in the
manifest. The all-size conclusion rests on
Lemmas~\ref{lem:overlap}--\ref{lem:last-active} and
Theorem~\ref{thm:abstract}, with the two independent exchange proofs.

The tests do, however, target likely proof errors: an incorrect exchange
direction, reversing one of the lexicographic orders, using the wrong active
set, missing an old face shared through a third base, forgetting the terminal
ridge, or confusing the ordinary and flag $h$-transforms. The independent
supply-versus-inequality comparison also checks the orientation of the
preorder in the matching model.

\subsection{How to reproduce the package}

From the extracted archive directory, run
\begin{verbatim}
python3 code/verify.py --output results
pdflatex -interaction=nonstopmode -halt-on-error article.tex
pdflatex -interaction=nonstopmode -halt-on-error article.tex
pdflatex -interaction=nonstopmode -halt-on-error article.tex
\end{verbatim}
The supplied \texttt{Makefile} provides the targets \texttt{verify},
\texttt{pdf}, and \texttt{clean}. Run Python without \texttt{-O}, because
assertions are the verifier's checks. The recorded JSON, CSV, and text output
were generated by the included code, not transcribed from a conjectural table.

The default exhaustive range is $n\le4$. A \texttt{--max-n 5} option is
provided for a larger run, but no size-five exhaustive run is included or
claimed. The face-complex verifier is intentionally exponential in rank: it
checks all $2^r$ subfaces of every facet. This is a diagnostic tool, not a
polynomial-time recognition algorithm for succinctly specified posets.

\section{Algorithmic content and proof audit}
\label{sec:audit}

\subsection{Generating the shelling without constructing the whole complex}

Given an explicit equidegree base list satisfying exchange, the following
procedure emits a shelling certificate:
\begin{enumerate}
\item Store the bases in a membership structure and sort them decreasingly
lexicographically.
\item For each base $b$, test $b+\ee_i-\ee_j\in\B$ for $i<j$ and $b_j>0$
to determine $R(b)$. Sort the alphabet as in \eqref{eq:alphabet}.
\item Generate distinct multiset words of content $b$ in that alphabet order.
For each word, emit its prefix chain and the ranks in \eqref{eq:restriction-ranks}.
\end{enumerate}
No repeated occurrence of an equal letter is distinguished; the word
recursion decrements remaining multiplicities rather than generating and
deduplicating permutations of labeled copies.

If $m=|\B|$, $s$ is the number of coordinates, and $N=N(D)$, there are at most
$m\binom{s}{2}$ transfer-membership tests, a sorting phase on $m$ vectors,
and $N$ emitted chains of length $r$. In a straightforward implementation
that handles an $s$-coordinate vector in $O(s)$ time, a conservative bound is
$O(ms^3+ms\log m+Nrs)$, assuming expected constant-time table access once a
key has been hashed. Integer bit costs and generation of the input base
list are not included in this bound. The output itself can be exponentially
large. The theorem gives a constructive order; it does not promise a small
listing when the complex has very many facets.

\subsection{Dependency audit}

The logical path can be read without any enumeration result from the
manifest:
\[
\begin{gathered}
 \text{ideal constraints}
 \Longrightarrow \text{supply assignments}
 \Longrightarrow \text{equidegree bases and exchange}\\
 \Longrightarrow \text{codimension-one box overlaps}
 \Longrightarrow \text{adapted word order}\\
 \Longrightarrow \text{exact old-ridge identity}
 \Longrightarrow \text{ordinary shellability}.
\end{gathered}
\]
The alternative input path replaces supply assignments by a coverage rank
function and the tight-set proof. The Cohen--Macaulay conclusion uses only
the standard homology consequence of a pure shelling, whose relevant argument
was included. The support retraction is independent of the exchange proof.
The optional connection to an Ehrhart $h^*$-polynomial is the only place
where an enumerative theorem of \cite{AC} is imported.

A particularly important audit point is that an old \emph{chain face} below
some earlier base must extend to a full chain in that base's block. Merely
showing that its individual vertices are in the union of earlier boxes would
not suffice in an arbitrary covering. Here its largest vertex is below one
earlier base; then the entire face is in that same box. This is why
Lemma~\ref{lem:last-active} proves an identity of complexes, not only an
identity of vertex sets.

Another audit point is the orientation in the exchange lemma. We specify the
coordinate to be \emph{increased}, because the first lexicographic difference
has $c_i>b_i$. The surplus coordinate to be decreased must occur later than
$i$. Neither a reversed exchange statement nor a generic assertion of
connectedness of the base graph would establish \eqref{eq:overlap}.

\subsection{Limits of the claim and the remaining questions}

The result answers the shellability and Cohen--Macaulayness question. It does
not supply an EL-labeling of the augmented poset: the letter order depends
on the terminal base and need not assign one consistent label to a cover
relation. No EL-shellability assertion is made. It also does not settle the
source's stronger real-rootedness, full $q$-reciprocity, or flag-polytope
questions \cite{AC}.

Two natural next investigations are whether the explicit order can be
realized by a suitable chain-labeling scheme, and whether the restriction
statistics admit simpler formulas independent of base-dependent alphabets.
These are follow-on questions, not missing hypotheses in the present
shellability proof. The original target already follows from the exact
intersection calculation.

\section{Conclusion}

The missing gluing step is controlled by unit exchange. It says that an
earlier base can meet a new coordinate box only through a codimension-one
box obtained by decrementing an active coordinate. Putting those coordinates
last in the local alphabet makes the final active occurrence either an
ordinary descent or the terminal step. Hence every old face lies in an
explicit old ridge, and the entire old intersection is exactly a union of
such ridges.

This proves the claimed affirmative answer to Question 4.6 for every finite
preorder, without the existence of a maximum vertex. It also gives a
self-contained shelling construction for integral-polymatroid integer-point
posets, together with their restriction faces and flag enumerators. The
included finite verification provides reproducible checks of the construction
and its conventions; the general proof remains the mathematical basis for
the result.

\begin{thebibliography}{9}

\bibitem{AC}
C.~A. Athanasiadis and F.~Chapoton,
\emph{Polytopes and posets associated to preorders},
arXiv:2605.26916v1 (2026), especially Question~4.6, p.~13.
\url{https://arxiv.org/abs/2605.26916v1}.

\bibitem{HH}
J.~Herzog and T.~Hibi,
\emph{Discrete polymatroids},
Journal of Algebraic Combinatorics \textbf{16} (2002), 239--268.
\url{https://doi.org/10.1023/A:1021852421716}.

\bibitem{AGY}
M.~Alizadeh, A.~Goodarzi, and S.~Yassemi,
\emph{$M$-shellability of discrete polymatroids},
arXiv:1012.1075v1 (2010).
\url{https://arxiv.org/abs/1012.1075v1}.

\bibitem{Ficarra}
A.~Ficarra,
\emph{Shellability of componentwise discrete polymatroids},
arXiv:2312.13006v2 (2023).
\url{https://arxiv.org/abs/2312.13006v2}.
Published in the Electronic Journal of Combinatorics (2025),
\url{https://doi.org/10.37236/12818}.

\bibitem{Wachs}
M.~L. Wachs,
\emph{Poset topology: tools and applications},
arXiv:math/0602226v2 (2006).
\url{https://arxiv.org/abs/math/0602226v2}.

\bibitem{DHLTW}
Z.~Dai, Q.~Hou, Z.~Liu, W.~Thawinrak, and H.~Wang,
\emph{Counting lattice points in Minkowski sums of cross polytopes},
arXiv:2608.16037v2 (2026), Section~4.3.
\url{https://arxiv.org/abs/2608.16037v2}.

\end{thebibliography}
\end{document}
